\begin{tikzpicture}[
    font=\small,
    % --- Stylings ---
    input/.style={
        draw=navy, fill=gray, rounded corners=3pt, align=center, 
        font=\footnotesize\bfseries, text=dark, inner sep=5pt, 
        minimum width=1.5cm, minimum height=0.75cm
    },
    param/.style={input, minimum width=0.8cm},
    model/.style={
        draw=navy, fill=white, rounded corners=4pt, thick, align=center, 
        font=\footnotesize\bfseries, text=navy, inner sep=4pt, 
        minimum width=1.4cm, minimum height=0.75cm
    },
    mathloss/.style={
        draw=teal, fill=teal!8, rounded corners=4pt, thick, align=center, 
        font=\small, text=dark, inner sep=6pt, 
        minimum height=1.5cm
    },
    weight/.style={
        draw=amber, fill=amber!15, circle, font=\scriptsize\bfseries, 
        text=dark, inner sep=1pt, minimum size=2.4em
    },
    sum/.style={
        draw=navy, fill=navy!15, circle, thick, font=\bfseries\large, 
        text=navy, inner sep=2pt
    },
    arr/.style={-Latex, thick, draw=navy},
    bus/.style={thick, draw=navy},
    tap/.style={circle, fill=navy, inner sep=0pt, minimum size=3.5pt},
    % NEU: Einheitlicher Style für Daten-Labels auf den Linien
    val/.style={
        anchor=south west, font=\tiny\bfseries, 
        inner sep=1pt, xshift=2pt, yshift=2pt
    }
]

    % ==========================================
    % 0. Y-GRID (Höhen-Schichten mit viel Platz)
    % ==========================================
    \coordinate (Y_sets)         at (0, 0);
    \coordinate (Y_split1)       at (0, -1.0); % Split für x
    
    \coordinate (Y_models1)      at (0, -2.2); % V_theta(x) & f(...)
    \coordinate (Y_split2)       at (0, -3.4); % Split für x_V^+
    
    \coordinate (Y_models2)      at (0, -4.6); % V_theta(x_V^+)
    
    \coordinate (Y_split3)       at (0, -6.0); % Routing V_theta(x)
    \coordinate (Y_split4)       at (0, -6.6); % Routing V_theta(x_V^+)
    
    \coordinate (Y_loss)         at (0, -9.2); % Mathe-Blöcke (benötigen Platz)
    
    \coordinate (Y_sum_bus)      at (0, -11.8);
    \coordinate (Y_sum)          at (0, -11.8);
    \coordinate (Y_lcond)        at (0, -13.2); % Combined Condition Loss
    
    \coordinate (Y_merge)        at (0, -14.8); % Routing für das Finale Gate
    \coordinate (Y_final)        at (0, -16.4); % Gated Loss

    % ==========================================
    % 1. X-GRID & EINGÄNGE 
    % ==========================================
    \node[param] (in_rho)   at (-5.5, 0 |- Y_sets) {$\rho$};
    \node[input] (in_x)     at (-0.5, 0 |- Y_sets) {$\mathcal{B}_\rho$};

    % ==========================================
    % 2. MODELLE & DATEN-SPLITS
    % ==========================================
    % Ebene 1
    \node[model] (net_v_x) at (-2.5, 0 |- Y_models1) {$V_\theta$};
    \node[model] (sys_f)   at ( 1.5, 0 |- Y_models1) {$f(x, \pi_V(x))$};
    
    \draw[arr] (in_x.south) -- (-0.5, 0 |- Y_split1) -| (net_v_x.north);
    \node[val] at (-2.5, 0 |- net_v_x.north) {$x$};
    
    \draw[arr] (-0.5, 0 |- Y_split1) -| (sys_f.north);
    \node[val] at (1.5, 0 |- sys_f.north) {$x$};

    % Ebene 2
    \node[model] (net_v_xp) at ( 1.5, 0 |- Y_models2) {$V_\theta$};
    
    % Routing für f(...) -> x^+
    \coordinate (tap_f) at (1.5, 0 |- Y_split2);
    \node[tap] at (tap_f) {};
    \draw[arr] (sys_f.south) -- (net_v_xp.north);
    \node[val] at (1.5, 0 |- net_v_xp.north) {$x^+$};

    % ==========================================
    % 3. MATH LOSS BOXEN (inkl. Formeln)
    % ==========================================
    
    % 3.1 Soft Sublevel Weight (w_rho)
    \node[mathloss] (w_rho) at (-5.0, 0 |- Y_lcond) {
        \textbf{Soft Sublevel Weight}\\[4pt]
        $\displaystyle w_{\rho} = \sigma \left( s \cdot \frac{\rho - V_\theta(x)}{\max (\rho, \varepsilon_{\text{rel}})} \right)$
    };

    % 3.2 Relative Decrease Loss (L_dec)
    \node[mathloss] (l_dec) at (-0.5, 0 |- Y_loss) {
        \textbf{Rel. Decrease Loss} $\mathcal{L}_{\text{dec}}$\\[4pt]
        $\displaystyle = \left[ \frac{(1 - \kappa) V_\theta(x) - V_\theta(x^+_V)}{\max(|V_\theta(x)|, \varepsilon_{\text{rel}})} \right]_+$
    };

    % 3.3 Relative Invariance Loss (L_inv)
    \node[mathloss] (l_inv) at (4.5, 0 |- Y_loss) {
        \textbf{Rel. Invariance Loss} $\mathcal{L}_{\text{inv}}$\\[4pt]
        $\begin{aligned}
            = & \left\| S_{\text{range}}^{-1} \big[ x^+_V - \bar{x}_V \big]_+ \right\|_1 \\[-2pt]
            + & \left\| S_{\text{range}}^{-1} \big[ \underline{x}_V - x^+_V \big]_+ \right\|_1
        \end{aligned}$
    };

    % ==========================================
    % 4. WATERFALL ROUTING ZU DEN LOSS-BOXEN
    % ==========================================
    \draw[arr] (-5.5, 0 |- in_rho.south) -- (-5.5, 0 |- w_rho.north);
    \node[val] at (-5.5, 0 |- w_rho.north) {$\rho$};

    % V_theta(x) Routing
    \coordinate (tap_vx) at (-2.5, 0 |- Y_split3);
    \draw[bus] (net_v_x.south) -- (tap_vx);
    \node[tap] at (tap_vx) {};
    
    % -> zu w_rho
    \draw[arr] (tap_vx) -| (-4.5, 0 |- w_rho.north);
    \node[val] at (-4.5, 0 |- w_rho.north) {$V_\theta(x)$};
    
    % -> zu L_dec
    \draw[arr] (tap_vx) -| (-1.0, 0 |- l_dec.north);
    \node[val] at (-1.0, 0 |- l_dec.north) {$V_\theta(x)$};

    % -> zu L_inv
    \draw[bus] (tap_f) -- (4.5, 0 |- Y_split2);
    \draw[arr] (4.5, 0 |- Y_split2) -- (4.5, 0 |- l_inv.north); % Korrigiert auf l_inv.north
    \node[val] at (4.5, 0 |- l_inv.north) {$x^+$};

    % V_theta(x_V^+) Routing
    \draw[arr] (net_v_xp.south) |- (0.0, 0 |- Y_split4) -- (0.0, 0 |- l_dec.north);
    \node[val] at (0.0, 0 |- l_dec.north) {$V_\theta(x_V^+)$};


    % ==========================================
    % 5. AGGREGATION: COMBINED CONDITION LOSS
    % ==========================================
    \node[sum] (sum_cond) at (2.0, 0 |- Y_sum) {$+$};
    
    % Routing L_dec -> Sum
    \draw[arr] (l_dec.south) |- (sum_cond.west);
    
    % Routing L_inv -> Sum (mit Gewicht)
    \draw[bus] (l_inv.south) -- (4.5, 0 |- Y_sum_bus);
    \node[weight] (w_inv) at (3.25, 0 |- Y_sum_bus) {$\omega_{\text{inv}}$};
    \draw[arr] (4.5, 0 |- Y_sum_bus) -- (w_inv.east);
    \draw[arr] (w_inv.west) -- (sum_cond.east);

    % Combined Loss Box
    \node[mathloss, minimum height=1.0cm] (l_cond) at (2.0, 0 |- Y_lcond) {
        \textbf{Combined Condition Loss}\\[4pt]
        $\displaystyle \mathcal{L}_{\text{cond}} = \mathcal{L}_{\text{dec}} + \omega_{\text{inv}} \cdot \mathcal{L}_{\text{inv}}$
    };
    \draw[arr] (sum_cond.south) -- (l_cond.north);


    % ==========================================
    % 6. FINALES GATE: L_rho-gated
    % ==========================================
    \node[mathloss, minimum height=1.2cm] (l_gated) at (-1.5, 0 |- Y_final) {
        \textbf{Combined $\rho$-Gated Condition Loss}\\[4pt]
        $\displaystyle \mathcal{L}_{\rho\text{-gated}} = \frac{ \sum_{\mathcal{B}_{\rho}} w_{\rho} \mathcal{L}_{\text{cond}} }{ \sum_{\mathcal{B}_{\rho}} w_{\rho} }$
    };

    % Zusammenführung von w_rho und L_cond
    \draw[arr] (w_rho.south) |- (-2.5, 0 |- Y_merge) -- (-2.5, 0 |- l_gated.north);
    \node[val] at (-2.5, 0 |- l_gated.north) {$w_\rho$};
    
    \draw[arr] (l_cond.south) |- (-0.5, 0 |- Y_merge) -- (-0.5, 0 |- l_gated.north);
    \node[val] at (-0.5, 0 |- l_gated.north) {$\mathcal{L}_{\text{cond}}$};

\end{tikzpicture}